\section{Related Work}
\label{sec:related-work}

Citation-based similarity has long been studied through co-citation
\citep{small1973cocitation} and bibliographic coupling \citep{kessler1963bibliographic}.
These methods rely on shared neighborhood structure and remain strong baselines because
they are simple, interpretable, and effective in sparse scholarly graphs.

Path-based and recursive approaches extend these ideas by considering longer-range
dependencies in the network. Random-walk formulations such as PageRank
\citep{page1999pagerank} and structural similarity measures such as SimRank
\citep{jeh2002simrank} exemplify this direction. More recently, graph neural networks have
provided representation-learning methods that can incorporate citation structure together
with textual or metadata features \citep{kipf2017gcn}.

At the researcher level, author similarity is often derived from publication overlap,
co-authorship, topic distributions, or aggregated article embeddings. A central challenge
is preserving meaningful article-level signals while controlling for author productivity
and field-specific citation practices.
